% !PRINTED 608
\weberpage{608}
Now by \S.\,164, (8) the difference $\Theta - \Theta'$ is equal to the square root of the fundamental number of $\Omega$, so that from (6) and (8) we obtain:
\weq{9}{\omega = \frac{b+\sqrt{\Delta}}{2c}, \quad \Delta = b^2 + 4ac,}
wherein the sign of $\sqrt{\Delta}$ is determined by the assumption made in \S.\,164, (6) about the sign of $\sqrt{d}$.

Thus $\omega$ is a \gesperrt{quadratic irrational number} belonging to the discriminant $\Delta$, such as we have considered in \S.\,122 of the first Band, and the form $\varphi$ can, if we set $a_{2,2}=e$, be expressed thus:
\weq{10}{\varphi = ec(t_1+\omega t_2).}

Herein $c$ is the least positive integral rational number for which the product $c\omega$ becomes an integer, and every other integral rational number $k$ by which the product $k\omega$ becomes an integer must be divisible by $c$.

One recognises this most simply if, by \S.\,164, one puts $\omega$ in one of the two forms, according as $b$ is even or odd:
\weq{11}{\omega = \frac{1}{c}\left(\frac{b}{2}+\Theta\right), \quad \frac{1}{c}\left(\frac{b-1}{2}+\Theta\right),}
from which, since $1, \Theta$ is a basis of $\mathfrak{o}$, what has been said follows.

On this is now founded the proof that, if $\omega$ denotes any quadratic irrational number belonging to the discriminant $\Delta$, $c$ the least positive integral rational number that makes $c\omega$ an integer, and $e$ an arbitrary positive integral rational number, then conversely, by
\weq{12}{\varphi = ec(t_1+\omega t_2)}
a basis form is always represented.

For this we have only to show that
\[
\alpha_1 = ec, \quad \alpha_2 = ec\omega
\]
is a basis of the functional $\varphi$ defined by (12).

If for this purpose we first form the norm of $\varphi$, we obtain by (7)
\weq{13}{N(\varphi) = e^2 c(ct_1^2 + bt_1t_2 - at_2^2),}
and consequently the absolute norm
\weq{14}{N_a(\varphi) = e^2 c.}
